\documentclass[11pt,a4paper]{article}
\usepackage[T1]{fontenc}
\usepackage{lmodern,amsmath,amssymb,amsthm,mathtools,microtype}
\usepackage[margin=27mm]{geometry}
\usepackage[colorlinks=true,linkcolor=blue,citecolor=blue,urlcolor=blue]{hyperref}
\setlength{\emergencystretch}{2em}
\newtheorem{proposition}{Proposition}[section]
\newtheorem{lemma}[proposition]{Lemma}
\newcommand{\Q}{\mathbb Q}\newcommand{\Z}{\mathbb Z}
\DeclareMathOperator{\Tr}{Tr}\DeclareMathOperator{\Nm}{N}
\title{An unshifted norm-square detector}
\author{}\date{17 November 2026}
\begin{document}\maketitle
\begin{abstract}
The single square condition $w\Nm(1-W^{-1})\in\Q^{\times2}$ gives
local elliptic index extraction for the fast conditional detector.
On actual multiplication tuples, its rational-square condition can be
prepared at every fresh split contact by one additional congruence on
the auxiliary index. The prescribed leading ratio is the integer input,
and requires no varying twist class. This proves soundness and fresh
contact preparation, not global rational realization.
\end{abstract}

\section{The square condition}

Retain the curve, point charts, six symmetric tests and contact function
of \cite{Single}. Thus
\[
 \begin{gathered}
 K=\Q(\sqrt2),\quad E:y^2=x^3-Dx,\quad D=(3-\sqrt2)^2,
 \quad 2E(K)=\Z P,\\
 F(t)=\frac{t(t+1)(2t+1)(3t+1)(4t+1)(5t+2)}{120},
 \qquad w=\frac{F(t)}{4(t-1)^8},\\
 W=\prod_{i=1}^6\left(1+\frac{C_i^2}{F(t)}\right),
 \quad V=\Tr(W^{-1}),\quad J=\Nm(W^{-1}).
 \end{gathered}
\]
Keep every point, difference and rational-function guard, and consider
the single condition
\begin{equation}\label{eq:R0}
                 {\cal R}_0=w(1-V+J)
                    =w\Nm(1-W^{-1})\in\Q^{\times2}.
\end{equation}

\begin{proposition}\label{prop:extract}
At every nonexceptional odd contact of $F(t)$, a guarded solution of
\eqref{eq:R0} has a place above that contact where all six tests
have positive valuation. More precisely, if
$\nu=v_p(F(t))>0$ is odd, then
$v_{\mathfrak p}(C_i)>\nu/2$ for all $i$ at some place
$\mathfrak p\mid p$.
\end{proposition}
\begin{proof}
The contact normalizer is a unit, so $v_p(w)=\nu$.
At each place above $p$, every factor of $W$ has valuation
\[
        \min\{0,2v_{\mathfrak p}(C_i)-\nu\}\le0.
\]
The two candidate valuations cannot tie. Its inverse is integral,
with residue one when the factor is a unit and zero otherwise.
Consequently each local value of $W^{-1}$ has residue zero or one.

If all of these residues were zero, then
$\Nm(1-W^{-1})\equiv1\pmod p$. This would make
$v_p({\cal R}_0)=\nu$ odd, contrary to \eqref{eq:R0}.
Thus at least one local value of $W$ is a unit. Every summand in
its valuation is nonpositive, so each must be zero. Therefore
$2v_{\mathfrak p}(C_i)-\nu>0$ for every $i$.
\end{proof}

The original formal comparison then gives $X\equiv n\pmod p$
when $A=[n]P$. The mixed-power height argument of \cite{Single}
is unchanged. Thus its sound nine-coordinate criterion remains
sound with \eqref{eq:R0} in place of its larger radical.
This statement retains the source's subgroup, trace-height and
weak-$abc$ hypotheses.

At a prepared real tuple $t>1$, each real detector slot exceeds one:
the factor containing $f(H)$ is strictly greater than one because
the source guards exclude $f(H)=0$. Hence
\[
            w>0,\qquad \Nm(1-W^{-1})>0.
\]
The same argument holds at every geometric point-algebra character.
The final radicand is therefore a unit at every prepared character.
One quadratic root gives rank $128$ and the same nine-coordinate
census, conditional on actual rational realization.

\section{Square preparation at a fresh contact}

Fix an integer input $X=n\ne0$, put $A=[n]P$, and restrict to
the actual multiplication family
\[
                        H=[h]P,\qquad U=[n]H.
\]
Let $p$ be a split good contact prime outside the fixed exceptional
set and the support of $n(1+n^4)$, with $v_p(F(t))=1$.
Choose the successful slot so that $H$ is formal there and $H-P$
is formal at the conjugate slot. Write
\[
 \tau=\zeta(H),\qquad
 \tau'=\overline{\zeta(H-P)},\qquad \zeta=-x/y.
\]

\begin{lemma}\label{lem:lead}
If $v_p(\tau)=v_p(\tau')=m\ge1$ and
$\tau/\tau'\equiv\varepsilon n\pmod p$ for
$\varepsilon\in\{1,-1\}$, then ${\cal R}_0$ is a nonzero
square in $\Q_p$.
\end{lemma}
\begin{proof}
The integral formal series give
\[
 f(Q)=\zeta(Q)^2(1+O(p^{4m})),\qquad
 \frac{\zeta([n]Q)}{\zeta(Q)}=n+O(p^{4m}).
\]
The first four tests therefore give
\[
 C_1^2+C_2^2+C_3^2+C_4^2
 =(1+n^4)(\tau^4+(\tau')^4)+O(p^{8m}).
\]
Both comparison tests have valuation at least $4m$, so the same
leading expression holds for $\sum_{i=1}^6C_i^2$.
The prescribed ratio makes this expression congruent to
\[
                         (1+n^4)^2(\tau')^4
                           \pmod {p^{4m+1}},
\]
with nonzero leading coefficient.

At the successful slot,
\[
 W-1=\frac{\sum_iC_i^2}{F(t)}(1+O(p)).
\]
At the other slot, the mandatory reciprocal test of $H$ is a unit:
there $H$ reduces to $P$, whose nonzero abscissa is a unit outside
the source's fixed exceptional set. Thus $W^{-1}=O(p)$ at that
slot. Taking the norm and using $w/F(t)=1/(4(t-1)^8)$ gives
\[
 {\cal R}_0=\frac{\sum_iC_i^2}{4(t-1)^8}(1+O(p)).
\]
Its valuation is $4m$ and its unit part is a square modulo $p$.
Hensel's lemma proves the assertion.
\end{proof}

\begin{proposition}\label{prop:chooseh}
The additional leading-ratio condition in Lemma~\ref{lem:lead}
can be imposed by a $p$-power congruence on $h$ compatible with
the original conjugate reduction congruences, provided $p$ is
coprime to their total modulus.
\end{proposition}
\begin{proof}
For the two embeddings define the formal logarithms of $P$ by
\[
 L_i=r_i^{-1}\log_E([r_i]P),
\]
where $r_i$ is a prime-to-$p$ annihilating reduction order. These
are nonzero elements of $p\Z_p$. On the original reduction
progression, the relevant logarithms are $hL_1$ and $(h-1)L_2$.

Choose $\varepsilon\in\{1,-1\}$ so that
\[
 v_p(L_1-\varepsilon nL_2)
                   =\min\{v_p(L_1),v_p(L_2)\}.
\]
If the valuations differ either sign works. If they agree, at
least one sign works because $p\ne2$ and $n$ is a unit. Then
\[
 h_*=-\frac{\varepsilon nL_2}{L_1-\varepsilon nL_2},\qquad
 h_*-1=-\frac{L_1}{L_1-\varepsilon nL_2}
\]
are nonzero elements of $\Z_p$. The logarithms have common
valuation $m=\max\{v_p(L_1),v_p(L_2)\}$ and exact ratio
$\varepsilon n$. The formal exponential differs from its linear
term by relative order at least $4m$, so their formal parameters
have the required leading ratio. A sufficiently precise congruence
$h\equiv h_*\pmod {p^k}$ preserves this property. The Chinese
remainder theorem combines it with the original modulus.
\end{proof}

\section{Simultaneous compatibility and exact verification}

For the original unrestricted rational-contact source, the
prime-pattern theorem chooses contacts of the form
$u+c_iw_1^2$, with fixed integer $c_i$. By increasing its prescribed
real lower bound on $u/w_1^2$, all fresh primes can be taken in an
interval $[P_0,2P_0]$ with $P_0$ arbitrarily large. Each split
reduction has full rational two-torsion, so its group order is a
multiple of four. Hasse's bound makes that order less than $4P_0$.
No fresh prime can therefore divide any of these group orders.
The finite initial modulus and the support of $n(1+n^4)$ are
avoided in advance. Thus all additional $p$-power congruences in
Proposition~\ref{prop:chooseh} are simultaneously compatible with
the original order progression. The retained real-orbit argument
continues to apply after this finite congruence refinement.
This use of the prime-pattern theorem does not impose the additional
thin condition $t=(\beta^e+T^e)^d+z^2$. For the nine-coordinate
criterion, the local lemmas apply to any contact already satisfying
their hypotheses; existence of a suitable thin argument remains part
of the unproved rational realization problem.

The checker accompanying this manuscript independently computes
the formal logarithm units at $p=77017$ using Jacobian arithmetic
modulo $p^2$. The conjugate annihilating orders are $2141,2405$;
the logarithm units are $8412,71610$. It verifies the leading-square
identity and compatible index congruences for inputs
$1,2,3,5,8,-2,-5$. For example, at $n=2$ it obtains
$h\equiv7796\pmod {77017}$, with formal leading units
$38485,57751$ and squared leading root $36042$ modulo $77017$.
The full combined congruence is recorded exactly in the certificate.

These results do not provide the required rational square over
$\Q$. Preparation at all fixed places is a separate task, and
fresh-contact local squares cannot exclude new odd valuations
arising from the rational functions $\Nm(W)$ and $\Nm(W-1)$.
In particular no global factorization on the actual multiplication
family, no complete rational realization theorem, and no improved
defining-variable bound is asserted.

\begin{thebibliography}{9}\raggedright
\bibitem{Single}\emph{A single-cover criterion from a mixed-power height
inequality}, repository manuscript,
\nolinkurl{papers/pole_parameters/single_cover_criterion_20261117_v56.tex}.
\end{thebibliography}
\end{document}
